% !TeX program = xelatex
% ============================================================================
%  第 10 课　生成元、循环群与循环群的子群　—— 教师版讲义（原 12+13 合并）
%  与 02-课件/第10课-生成元与循环群/第10课-生成元与循环群.tex 同步
% ============================================================================

\jhlesson{10}{生成元、循环群与循环群的子群}{时长 45 分钟\quad·\quad 教具：12 格圆盘（在黑板上画；第 2 课画过一个）\quad·\quad 本课不发单页}

\begin{mubiao}
  \begin{itemize}
    \item 知道什么是生成元，能说出 \(D_4\) 里哪两个就够用。
    \item 能找出模 \(12\) 加法里由 \(3\)、由 \(5\) 生成的子群。
    \item 会求模 \(12\) 各元素的阶，并发现子群大小都是 \(12\) 的因数。
  \end{itemize}
\end{mubiao}

\noindent{\small 本课节奏：0—6 回顾｜6—20 生成元｜20—32 各元素的阶｜32—43 循环群的子群｜43—45 收尾。}

\section*{一、回顾（0—6 分钟）}

上节课最后留了一句：\(D_4\) 的 \(8\) 个动作，只要 \(r\) 和 \(f\) 就够。今天先回答「够」是什么意思。

板书把其余六个用 \(r\)、\(f\) 拼出来：\(r^2=rr\)，\(r^3=rrr\)，\(rf\)，\(r^2f\)，\(r^3f\)。

\section*{二、生成元（6—20 分钟）}

\begin{definition}[由 \(a\) 生成的子群]
  在群 \(G\) 里，把元素 \(a\) 和它反复做出来的所有结果放在一起，
  得到的那个小群，叫做\textbf{由 \(a\) 生成的子群}，记作 \(\langle a\rangle\)。
\end{definition}

\begin{definition}[生成元、循环群]
  如果整个群 \(G\) 能由一个元素 \(a\) 生成，也就是 \(G=\langle a\rangle\)，
  就说 \(a\) 是 \(G\) 的一个\textbf{生成元}，而 \(G\) 是一个\textbf{循环群}。
\end{definition}

\begin{caozuo}
  在黑板上画一个 12 格圆盘，当着全班走两遍：
  \begin{itemize}
    \item 从 \(0\) 开始一直加 \(3\)：\(0\to3\to6\to9\to12\to\cdots\)，
      走到 \(12\) 就等于回到 \(0\)，所以 \(\langle3\rangle=\{0,3,6,9\}\)，一共 \(4\) 个；
    \item 从 \(0\) 开始一直加 \(5\)：\(0\to5\to10\to3\to8\to1\to6\to11\to4\to9\to2\to7\to0\)，
      把 \(12\) 个数一个不落全走了一遍，所以 \(5\) 是生成元，这个群是循环群。
  \end{itemize}
\end{caozuo}

然后抛问题：为什么 \(3\) 只圈住 \(4\) 个，\(5\) 却圈住 \(12\) 个？

\begin{center}
\begin{tabular}{@{}lcc@{}}
  \toprule
   & 和 \(12\) 的公因数 & 圈住几个 \\
  \midrule
  \(3\) & \(1,3\)（最大的是 \(3\)） & \(12\div3=4\) \\
  \(5\) & 只有 \(1\) & \(12\div1=12\) \\
  \bottomrule
\end{tabular}
\end{center}

一句话：公因数越大，圈住的越少；只有公因数 \(1\) 时，整整一圈都被圈住。

\begin{zhu}
  这里可以和第 3 课对一下：第 3 课讲的是模 \(7\) 的乘法（人人有逆元）和模 \(6\) 的乘法（\(2\) 没有）；
  换成模 \(12\) 也一样——和 \(12\) 只有公因数 \(1\) 的数，才有逆元。
  两件事是一回事，不必展开证明。
\end{zhu}

\begin{caozuo}
  手指回答：\(\langle2\rangle\) 里有几个数、是哪几个？（\(6\) 个：\(\{0,2,4,6,8,10\}\)）
  \(\langle6\rangle\) 呢？（\(2\) 个：\(\{0,6\}\)）\(7\) 是生成元吗？（是，\(7\) 与 \(12\) 只有公因数 \(1\)）
\end{caozuo}

\section*{三、各元素的阶（20—32 分钟）}

板书一张空表，把 \(1\text{—}11\) 和 \(0\) 写在上面。全班用手指数报「几次回到 \(0\)」，你只管填。

\begin{center}
\begin{tabular}{@{}l*{12}{c}@{}}
  \toprule
  \(a\) & 1 & 2 & 3 & 4 & 5 & 6 & 7 & 8 & 9 & 10 & 11 & 0 \\
  \midrule
  阶 & 12 & 6 & 4 & 3 & 12 & 2 & 12 & 3 & 4 & 6 & 12 & 1 \\
  \bottomrule
\end{tabular}
\end{center}

填完点一句：这就是第 7 课那个「做几次回到原地」，只不过这次做的是数。

\begin{dingli}[阶与生成子群]
  元素 \(a\) 的阶，正好等于由它生成的子群 \(\langle a\rangle\) 的大小。
\end{dingli}

\begin{caozuo}
  手指回答：\(4\) 的阶是多少？\(\langle4\rangle\) 里有几个数？（都是 \(3\)，\(\{0,4,8\}\)）
  \(8\) 呢？\(\langle8\rangle\) 和 \(\langle4\rangle\) 是同一个子群吗？（是，\(\{0,8,4\}\)）
\end{caozuo}

\begin{zhu}
  学生版那道「\(\langle4\rangle\) 和 \(\langle8\rangle\) 是同一个子群吗」的答案口径：\textbf{是}，
  \(\langle4\rangle=\langle8\rangle=\{0,4,8\}\)——因为它们有相同的公因数
  （都是 \(4\) 与 \(12\) 的最大公因数）。
\end{zhu}

\section*{四、循环群的子群（32—43 分钟）}

在 12 个数的圈上，把所有「能自成一体的小圈」一个一个圈出来：

\begin{center}
\begin{tabular}{@{}clc@{}}
  \toprule
  子群 & 是谁 & 大小 \\
  \midrule
  \(\{0\}\) & 只有 \(0\) & \(1\) \\
  \(\{0,6\}\) & 加 \(6\) 来回走 & \(2\) \\
  \(\{0,4,8\}\) & 由 \(4\) 生成 & \(3\) \\
  \(\{0,3,6,9\}\) & 由 \(3\) 生成 & \(4\) \\
  \(\{0,2,4,6,8,10\}\) & 由 \(2\) 生成 & \(6\) \\
  \(\{0,1,2,\dots,11\}\) & 全部十二个 & \(12\) \\
  \bottomrule
\end{tabular}
\end{center}

一共就这 \(6\) 个，没有别的。数它们的大小：\(1,2,3,4,6,12\)——每一个都能整除 \(12\)。

\begin{zhu}
  这句话只留印象，不要求证明。三张纸对起来看：
  \(S_3\) 里是 \(1,2,3,6\) 整除 \(6\)；\(D_4\) 里是 \(1,2,4,8\) 整除 \(8\)；
  这里是 \(1,2,3,4,6,12\) 整除 \(12\)。同一句话，第三个例子。
\end{zhu}

\section*{五、收尾（43—45 分钟）}

板书三句话：一个元素反复做能造出整个群，就叫生成元，这样的群叫循环群；
元素的阶等于它生成的子群的大小；循环群的子群大小一定整除整个群的大小。

\begin{xiaojie}
  \begin{itemize}
    \item \textbf{生成元}：\(\langle3\rangle=\{0,3,6,9\}\)，\(\langle5\rangle\) 是全部的 \(12\) 个；模 \(12\) 加法是循环群。
    \item \textbf{阶 \(=\) 生成子群的大小}：\(a\) 的阶就是 \(\langle a\rangle\) 里元素的个数。
    \item 循环群的子群大小 \(1,2,3,4,6,12\) 都能整除 \(12\)。
  \end{itemize}
\end{xiaojie}

\begin{zhu}
  \textbf{本课常见错误，提前打预防针：}
  \begin{itemize}
    \item 把「生成元」理解成「群里面最大的元素」——不是，是「一个就够用」的那个元素；
    \item 把 \(\langle 3\rangle\) 与 \(\langle 9\rangle\) 当成两个不同的子群——它们其实是同一个
      （\(\{0,3,6,9\}\)），换个人当头而已；
    \item 以为 \(\langle 5\rangle\) 只有 \(5\) 个元素——它走遍全部 \(12\) 个。
  \end{itemize}
\end{zhu}

\begin{zhu}
  \textbf{时间不够就砍这里：}把「\(\langle 4\rangle\) 与 \(\langle 8\rangle\) 是同一个子群」那一段
  改成老师直接报结论，省下的 \(2\) 分钟留给最后的「三张纸对起来看」。
  最后一页（子群大小整除群的大小）是第 11、12 课的地基，不能砍。
\end{zhu}

\noindent\textbf{下节课}：回头去看幂的尾巴——\(2\) 的幂除以 \(7\)，余数会转圈。